
% Annexes
\part{Annexes}
\section{exclusion des autres caractères}
\label{sec:exclu}
\noindent

\begin{tabular}{p{0.33\textwidth}p{0.66\textwidth}}
  Caractère(s) & Raison\\ 
  \hline
  Caractère hors du jeu ASCII 7~bits 
  et l'espace.
  & Même si les
  systèmes de fichier des systèmes d’exploitation actuels savent gérer accents
  , espaces et caractères non-latins, il subsiste des problèmes de portabilité
  (notamment lors de l'utilisation de NFS ou SMB entre systèmes avec 
  des codages différents), et il n'y
  a pas suffisamment d'avantages à utiliser les accents pour accepter les
  inconvénients associés.\\
  \hline
  Les caractères de contrôle, de code entre 0 et
  31, et 127(« del ») & Ne sont pas faits pour cet usage. \\
  \hline
  \texttt{/} (barre oblique) et\newline
  \texttt{\textbackslash} (barre oblique inversée) &
   Ils servent de séparateurs de nom de répertoire dans les
  chemins d'accès aux fichiers sur différents systèmes d'exploitation et dans
  les URL/URI.\\
  \hline
  \texttt{"} (guillemet droit),\newline
  \texttt{*} (astérisque ou étoile),\newline
  \texttt{?} (point d'interrogation),\newline
  \texttt{>} et \texttt{<} (supérieur et inférieur),\newline
  \texttt{|} (barre verticale)
   & Ils sont impossibles sur certains systèmes de fichiers.\\
  \hline
  \texttt{'} (apostrophe),\newline
  \texttt{\$} (dollar),\newline
  \texttt{;} (point-virgule),\newline 
  \texttt{\&} (esperluette) et\newline
  \texttt{`} (guillemet ouvrant simple) & Les
  inconvénients qu'ils engendrent pour les utiliser correctement dans des
  environnements usuels (shell unix, par exemple) ne sont pas supérieurs à
  l'intérêt que pourrait avoir leur utilisation\\
  \hline
  \texttt{+} (plus),\newline
  \texttt{-} (moins, trait d'union ou tiret),\newline
  \texttt{\^{}} (accent circonflexe avec chasse),\newline
  \texttt{(} et \texttt{)} (parenthèses ouvrante et fermante),\newline
  \texttt{[} et \texttt{]} (crochets ouvrant et fermant) et\newline
  \texttt{=} (signe égal) & Ils servent aux opérations
  mathématiques usuelles lors de l'interprétation d'une formule (avec « * » et
  « / ». qui sont déjà impossibles pour d'autres raisons).\\
  \hline
\end{tabular}
